\section{Introduction}
\label{sec:introduction}

Given the shadow of a knot diagram---the four-valent planar projection with
crossing signs forgotten---two players alternately decide each crossing until
a complete diagram remains. One player, the \emph{knotter}, aims for a knotted
diagram; the other, the \emph{unknotter}, aims for the unknot. We call this
``to knot or not to knot.'' The game is finite, perfect-information, and
zero-sum, and for small crossing numbers its full game tree fits comfortably in
memory: a seven-crossing shadow admits $3^7 = 2{,}187$ partial states, of which
2,059 are nonterminal and 128 are terminal. Every position is therefore exactly
solvable by minimax, which makes the game an unusual laboratory: a neural
search system can be trained by self-play, yet its every decision can be graded
against ground truth.

We pursue two goals. The first is engineering: solve as many finite knot-game
instances as possible with AlphaZero-style self-play~\cite{silver2017mastering}
and Monte Carlo tree search. The second is scientific: inspect the learned
policy, value, and internal representations for recurring combinatorial
principles, embed knots by strategic similarity rather than by classical
invariants alone, and compare the two geometries in search of conjectures. This
paper reports the foundation on which both goals rest: a validated game
implementation with exact ground truth, a PD-native architecture with certified
representational capacity, a versioned corpus with frozen evaluation splits,
and a first exploratory map of learned strategy space.

\subsection{Contributions}
\label{sec:contributions}

\begin{itemize}
    \item A validated crossing-resolution game on oriented PD shadows with
        exhaustive minimax tables. The fixed seven-crossing production shadow
        has 2,187 partial states (2,059 nonterminal), 128 terminal resolutions
        of which 70 satisfy the Jones$=1$ terminal rule, a forced win for the
        player to move's opponent (the unknotter under the production
        knotter-first roles), and 14 mutually optimal opening actions
        (Section~\ref{sec:exact}).
    \item A PD-native port-graph transformer that reaches exhaustive
        \texttt{SOLVED: YES} by exact supervision: 2,059/2,059 optimal
        top-policy actions with 99.70\% mean optimal mass and 2,059/2,059
        correct value signs on the fixed shadow (Bouchet job 23989965, epoch
        512), and 149,319/149,319 on both criteria for one shared
        variable-size model over the 35-diagram prime catalog from three
        through eight crossings (Bouchet job 24128608, three architectures).
        Both are representational-capacity results using exact labels, not
        self-play discovery or generalization (Section~\ref{sec:results}).
    \item A versioned manifest with frozen knot-type-disjoint train, validation,
        and test splits (\texttt{corpus\_manifest.py}; ROADMAP G1.1/G2 gate~3),
        structural plus determinism validation of all $21 \times 256$
        eight-crossing terminals
        (\texttt{validate\_8crossing\_terminals.py}), and a documented negative
        result: the end-to-end AlphaZero self-play run (job 23623670) reaches
        only 71.73\% optimal top-policy actions and does not solve the fixed
        shadow (Section~\ref{sec:validation}).
    \item A Lean formalization of PD signs, crossing changes, and the game
        itself, with the general flip-involution theorem and three certified
        solved games (trefoil, figure-eight, $5_2$), stating its oracle trust
        boundary explicitly (Section~\ref{sec:lean}).
    \item An exploratory activation-diffusion/PHATE pilot over the solved
        shared checkpoint (job 24228941) with a preregistered control checklist
        that must pass before any strategic-geometry claim
        (Section~\ref{sec:representations}).
\end{itemize}

Explicit non-claims: the AlphaZero self-play run did not solve the fixed
shadow; shared capacity used exact labels and is not generalization; the
Jones$=1$ terminal rule is the game's definition of a win, not a general
unknot-recognition theorem (see \texttt{AGENTS.md} sources of truth).
